% Section 16 Header
\section{Section 16 — Dynamic Programming}

\begin{frame}[c]{\empty}
    \begin{center}
        \Large\textbf{\secname}
    \end{center}
\end{frame}

% Slide 1: Concept
\begin{frame}{What is Dynamic Programming (DP)?}
  \begin{itemize}
    \item \textbf{The Core Idea:} Solve complex problems by breaking them down into simpler, overlapping subproblems and caching their results to avoid redundant calculations.
    \item \textbf{Key Properties:}
      \begin{itemize}
        \item \textbf{Overlapping Subproblems:} Solving the same subproblems repeatedly.
        \item \textbf{Optimal Substructure:} The optimal solution of the problem can be constructed from optimal solutions of its subproblems.
      \end{itemize}
    \item \textbf{Approaches:}
      \begin{itemize}
        \item \textbf{Top-Down (Memoization):} Recursive, start from target, solve subproblems and save values in a cache.
        \item \textbf{Bottom-Up (Tabulation):} Iterative, start from base cases, fill a table (array/matrix) up to target.
      \end{itemize}
  \end{itemize}
\end{frame}

\begin{frame}{Dynamic Programming Sub-patterns}
  Most interview DP problems fall into standard subclasses:
  \begin{itemize}
    \item **Fibonacci Numbers:** `dp[i] = dp[i-1] + dp[i-2]` (Climbing Stairs).
    \item **0/1 Knapsack:** Items can be chosen at most once. Decide to include or exclude an item: `dp[i][w] = max(exclude, include)`.
    \item **Unbounded Knapsack / Coin Change:** Items can be chosen multiple times.
    \item **Longest Common Subsequence (LCS) / Edit Distance:** String matching.
    \item **Longest Increasing Subsequence (LIS):** Decisions based on previous elements.
  \end{itemize}
\end{frame}

% Problem 1: Climbing Stairs
\begin{frame}[c]{\empty}
    \begin{center}
        \Large \textbf{Problem 1: \href{https://leetcode.com/problems/climbing-stairs/}{Climbing Stairs}}
    \end{center}
\end{frame}

% Problem 2: 0/1 Knapsack (Concept)
\begin{frame}[c]{\empty}
    \begin{center}
        \Large \textbf{Concept: 0/1 Knapsack}
    \end{center}
\end{frame}

% Problem 3: Coin Change
\begin{frame}[c]{\empty}
    \begin{center}
        \Large \textbf{Problem 3: \href{https://leetcode.com/problems/coin-change/}{Coin Change}}
    \end{center}
\end{frame}

% Problem 4: Coin Change 2
\begin{frame}[c]{\empty}
    \begin{center}
        \Large \textbf{Problem 4: \href{https://leetcode.com/problems/coin-change-2/}{Coin Change 2}}
    \end{center}
\end{frame}

% Problem 5: Longest Increasing Subsequence
\begin{frame}[c]{\empty}
    \begin{center}
        \Large \textbf{Problem 5: \href{https://leetcode.com/problems/longest-increasing-subsequence/}{Longest Increasing Subsequence}}
    \end{center}
\end{frame}

% Problem 6: Word Break
\begin{frame}[c]{\empty}
    \begin{center}
        \Large \textbf{Problem 6: \href{https://leetcode.com/problems/word-break/}{Word Break}}
    \end{center}
\end{frame}
